% Human Validation Rubric for RQ1 (Citation Judge Validation)
% Original rubric in Dutch with English translation

\begin{table}[htbp]
\centering
\caption[Human validation rubric for RQ1]{Human validation rubric for evaluating causal evidence quality in Citation Judge validation (RQ1). Original rubric in Dutch with English translation.}
\label{tab:human-validation-rubric-rq1}

\small
\begin{tabular}{p{0.95\textwidth}}
\toprule
\textbf{Stappen Causal Narrative Check} / \textit{Steps for Causal Narrative Check} \\
\midrule
1. Worden de variabelen of synoniemen genoemd en in de zelfde richting als in motivatie (A$\to$B en niet B$\to$A)? \\
\textit{Are the variables or synonyms named and in the same direction as in motivation (A$\to$B and not B$\to$A)?} \\[0.5em]

2. Is er een relatie statistisch significante (p $<$ 0,05) relatie? \\
\textit{Is there a statistically significant (p $<$ 0.05) relationship?} \\[0.5em]

3. Is de relatie causaal (A veroorzaakt B)? \\
\textit{Is the relationship causal (A causes B)?} \\[0.5em]

4. Als A$\to$C veroorzaakt via mediator B is het indirect, benoem in dat geval de mediator in annotatie. Wij zijn opzoek naar directe causale relaties. \\
\textit{If A$\to$C is caused via mediator B it is indirect; in that case name the mediator in annotation. We are looking for direct causal relationships.} \\[0.5em]

5. Kijken of de populatie uit de bron overeenkomt met de populatie van de stelling. Als dit anders of beperkt is moet dat benoemd worden in een aparte kolom. \\
\textit{Check if the population from the source matches the population of the statement. If different or limited, this must be noted in a separate column.} \\[0.5em]

6. Algemene check van kwaliteit bewijs: is het een longitudinale studie, wetenschappelijke studie of blogpost? \\
\textit{General quality check of evidence: is it a longitudinal study, scientific study, or blog post?} \\
\bottomrule
\end{tabular}

\vspace{1em}

\textbf{Classificatie} / \textit{Classification}

\vspace{0.5em}

\begin{tabular}{|>{\columncolor{green!70!black}\color{white}}p{2cm}|c|p{8cm}|}
\hline
\textbf{Kleur/Color} & \textbf{Score} & \textbf{Notitie / Note} \\
\hline
\cellcolor{green!70!black}\textcolor{white}{Dark Green} & 1.0 & 
Hoogste kwaliteit bewijs, causaal + wetenschappelijk, longitudinaal, RCT. Wanneer het duidelijk causaal bewijs hoeft de rest van de bronnen niet bekeken te worden. \\
& & \textit{Highest quality evidence, causal + scientific, longitudinal, RCT. When clear causal evidence exists, the rest of the sources need not be reviewed.} \\
\hline
\cellcolor{green!50} Light Green & 0.8 & 
Wel causaal verband maar geen wetenschappelijke studie, bijv.\ blog medische professionals. \\
& & \textit{Causal relationship but no scientific study, e.g., blog by medical professionals.} \\
\hline
\cellcolor{yellow!80} Yellow & 0.5 & 
Relatie maar geen causaal verband, cross-sectioneel onderzoek. \\
& & \textit{Relationship but no causal link, cross-sectional study.} \\
\hline
\cellcolor{red!70} Red & 0.0 & 
Geen toegang tot link, onjuiste variabelen, geen verband. \\
& & \textit{No access to link, incorrect variables, no relationship.} \\
\hline
\cellcolor{gray!50} Gray & 0.0 & 
Als de bron echt niet duidelijk is of er iets vreemds aan de hand is. Goed beschrijven wat er niet duidelijk is. \\
& & \textit{If the source is really unclear or something strange is going on. Describe well what is not clear.} \\
\hline
\end{tabular}

\end{table}



